\documentclass[11pt]{article}
\usepackage[margin=1in]{geometry}
\usepackage{amsmath,amssymb,amsthm}
\usepackage{booktabs}
\usepackage[colorlinks=true,linkcolor=blue,citecolor=blue,urlcolor=blue]{hyperref}

\newtheorem{theorem}{Theorem}
\newtheorem{lemma}[theorem]{Lemma}
\newtheorem{proposition}[theorem]{Proposition}
\newtheorem{corollary}[theorem]{Corollary}
\newtheorem{conjecture}[theorem]{Conjecture}
\newtheorem{definition}[theorem]{Definition}
\theoremstyle{remark}
\newtheorem{remark}[theorem]{Remark}

\newcommand{\Z}{\mathbb{Z}}
\newcommand{\F}{\mathbb{F}}
\newcommand{\Q}{\mathbb{Q}}
\DeclareMathOperator{\lcm}{lcm}
\DeclareMathOperator{\ord}{ord}

\title{What the Second Congruence Buys:\\ Quantified Evidence on
Popovych's Conjecture\\ \large (standalone note --- v0.1)}
\author{Dan Shumow \and Claude (Anthropic Fable~5)}
\date{October 22, 2026}

\begin{document}
\maketitle

\noindent\textbf{Audit revision (22 October 2026).}
The support lower bounds below are conditional, not consequences of the
collapse upper bounds alone. Special finite-field collapses are possible.
The tracked PDF is historical and has not been rebuilt from this revision.
See \texttt{doc/integration-review-status.md}.

\begin{abstract}
Agrawal's conjecture asserts that for a prime $r \nmid n$ the congruence
$(X-1)^n \equiv X^n - 1 \pmod{n, X^r - 1}$ forces $n$ prime or
$n^2 \equiv 1 \pmod r$; Popovych (2009) conjectured that adding the second
congruence $(X+2)^n\equiv X^n+2$ repairs the statement.
We compare support models for Frobenius-matching constructions. The first
section has the exact inversion identity
$(X-1)^{p^{d/2}-1}=-X^{-1}$ for even $d=\ord_r(p)$.
For the second section, a characteristic-zero ideal argument proves that
$(\zeta_r^\gamma+2)/(\zeta_r+2)$ is not a unit for $\gamma\ne1$.
This excludes a global root-of-unity identity, not exceptional finite-field
collapses or missing order factors.

Under explicit section-specific support-sharpness hypotheses at a fixed
smoothness threshold, the pair's containing-torus dimension is
$\min\{d:d\mid r-1,\ d\ge3\}$, versus the single-section model's minimum
$2$. At $r=5$ the pair model has dimension $4$, requiring smooth
$p-1$, $p+1$ and $p^2+1$. The lower smoothness requirements are
\emph{conditional}; upper order bounds alone do not prove them.
Numerical order checks and parameter-specific Dickman estimates are finite
evidence and heuristics, not a proof that every feasible harvest is empty.
We make no priority claim, proof of either conjecture, or general
algorithmic impossibility claim.
\end{abstract}

\section{Introduction}

For coprime $n, r$ write $T(a, n, r)$ for the congruence
$(X + a)^n \equiv X^n + a \pmod{n,\ X^r - 1}$. Agrawal's conjecture
\cite{AKS,BP01} states that $T(-1, n, r)$ with $r$ prime, $r \nmid n$,
forces $n$ prime or $n^2 \equiv 1 \pmod r$; it is open, and believed
false on the strength of the Lenstra--Pomerance (L--P) heuristic
construction \cite{LPremarks}. Popovych \cite{Popovych09} proved a
generalization of the L--P counterexample mechanism and proposed the
repaired conjecture:

\begin{conjecture}[Popovych]\label{conj:pop}
If $r$ is a prime not dividing $n$ and both $T(-1, n, r)$ and
$T(2, n, r)$ hold, then $n$ is prime or $n^2 \equiv 1 \pmod r$.
\end{conjecture}

Its evidence base is thin in both directions: a plausibility argument
that the pair $\{X - 1, X + 2\}$ generates a large unit subgroup
\cite{Popovych09,Popovych12}, and the Primaboinca exhaustive search
($10^{10}<n<10^{17}$, nothing found). This note provides a conditional
comparison of containing support products: global model minima $3$ and $2$,
and dimensions $4$ and $2$ at $r=5$. At some other fixed moduli the two
models tie; a larger global minimum is not a strict increase in every cell.

We consider the \emph{Frobenius-matching} family containing the original
L--P sufficient construction \cite{LPremarks}. In the $\Phi_r$ factor
modulo each prime $p\mid n$, matching holds if $n\equiv
p^{j} \pmod{\lcm(r, \ord(X + a))}$, where the order is taken in
$(\F_p[X]/\Phi_r(X))^\times$, because $(X + a)^{p^j} \equiv X^{p^j} + a$
exactly (Frobenius and Fermat). A construction assembles a squarefree
$n = p_1 \cdots p_k$ meeting these congruences at every $p_i$
simultaneously. Requiring those orders to be $B$-smooth is the particular
assembly design analyzed here, not a necessary condition for every
conceivable construction. Within that design, the
support is where the two sections differ. On the $X=1$ factor one must also
impose $(1+a)^n\equiv1+a\pmod p$; matching in the $\Phi_r$ factor alone does
not imply this. It is automatic for a squarefree Carmichael number, including
when the base is zero modulo a factor.

\section{The collapse for $X - 1$ and its absence for $X + 2$}

Fix an odd prime $r \ne 3$, a prime $p \nmid 3r(2^r + 1)$, and let
$d = \ord_r(p)$. Write $\rho_r(p)$ and $\rho'_r(p)$ for the
multiplicative orders of $X - 1$ and $X + 2$ in
$(\F_p[X]/\Phi_r(X))^\times$; both divide $p^d - 1 = \prod_{e \mid d}
\Phi_e(p)$, a factorization indexed by \emph{cyclotomic levels} $e$,
with $\Phi_e(p)$ the order of the $e$-th cyclotomic torus over $\F_p$,
of dimension $\varphi(e)$.

\begin{lemma}[inversion collapse for the first section]\label{lem:collapse}
If $d$ is even, then $(X-1)^{p^{d/2} - 1} = -X^{-1}$ in
$\F_p[X]/\Phi_r(X)$, and hence $\rho_r(p) \mid 2r(p^{d/2} - 1) =
2r \prod_{e \mid d/2} \Phi_e(p)$.
\end{lemma}

\begin{proof}
$(\Z/r)^\times$ is cyclic, so its unique involution is $-1$ and
$p^{d/2} \equiv -1 \pmod r$. In each field component, with $X \mapsto
\zeta$ a primitive $r$-th root of unity, iterated Frobenius gives
$(\zeta - 1)^{p^{d/2}} = \zeta^{-1} - 1 = -\zeta^{-1}(\zeta - 1)$;
divide by $\zeta - 1 \ne 0$. Finally $(-X^{-1})^{2r} = 1$ since
$X^r \equiv 1 \bmod \Phi_r$.
\end{proof}

For $r = 5$, $d = 4$ this is V\'a\v{n}a's $\rho_5(p) \mid 10(p^2 - 1)$
\cite{Vana09}, the fact that makes the L--P construction's cost the
\emph{double} condition ($p - 1$ and $p + 1$ smooth: dimension
$1 + 1 = 2$) rather than the full $p^4 - 1$. The collapse exists because
the ratio under inversion is torsion,
with root-of-unity cofactor. For the second section the global conjugate
ratio has a different property:

\begin{lemma}[trivial stabilizer]\label{lem:nonunit}
For every $\gamma \in (\Z/r)^\times$ with $\gamma \ne 1$, the element
$(\zeta^\gamma + 2)/(\zeta + 2)$ is not a unit of $\Z[\zeta_r]$ --- in
particular not a root of unity. There is therefore no characteristic-zero
torsion-cofactor identity for these conjugate ratios. This does not forbid
accidental finite-field relations of the form
$(X+2)^{p^e-1}\in\mu_{2r}$.
\end{lemma}

\begin{proof}
First, $N(\zeta + 2) = \Phi_r(-2) = (2^r + 1)/3$ (from $X^r - 1 =
(X - 1)\Phi_r(X)$ at $X = -2$), which is coprime to $r$: by Fermat
$2^r \equiv 2 \pmod r$, so $2^r + 1 \equiv 3$, and $\gcd(r, 3) = 1$.
Let $\mathfrak q$ be a prime of $\Z[\zeta_r]$ dividing $(\zeta + 2)$;
then $\zeta \equiv -2 \pmod{\mathfrak q}$, and the image of $-2$ in the
residue field has multiplicative order exactly $r$ (it is the image of
$\zeta$). If $\mathfrak q$ also divided $(\zeta^\gamma + 2)$ then
$\zeta^\gamma \equiv -2 \equiv \zeta$, so $\zeta^{\gamma - 1} \equiv 1
\pmod{\mathfrak q}$, forcing $r \mid \gamma - 1$. Hence for
$\gamma \ne 1$ the ratio has a pole at $\mathfrak q$ and is not a unit.
\end{proof}

\begin{remark}[a special reduction]
For $r=13$ and $p=5$, one has $5^2\not\equiv1\pmod{13}$ but
$(X+2)^{24}\in\mu_{26}$ in $\F_5[X]/\Phi_{13}(X)$.
This permitted prime illustrates why the global lemma is not a
no-exception statement about reductions.
\end{remark}

\begin{remark}[two structural asides]\label{rem:asides}
(i) The $X = 1$ component of $T(2, n, r)$ is the base-$3$ Fermat
condition $3^n \equiv 3 \pmod n$ --- automatic for Carmichael-type $n$,
but nontrivial, unlike the first section's $0 = 0$. (ii) Modulo $p = 3$
one has $-2 \equiv 1$, so the two sections \emph{coincide}: the pair
constrains only the divisors $p > 3$ of $n$. Finitely many $p$ (those
dividing $2^r + 1$) make $X + 2$ a zero divisor and are likewise
excluded.
\end{remark}

\section{The burden jump}

\begin{definition}[sharpness hypothesis]\label{def:sharp}
Fix $B\ge2rd$. For the chosen section $f$ and containing product $P_f(p)$,
assume that every prime $\ell>B$ dividing $P_f(p)$ also divides $\ord(f)$.
For $X-1$ take $P_f=p^{d'}-1$, where $d'=d/2$ for even $d$ and $d'=d$
otherwise; for $X+2$ take $P_f=p^d-1$. These are separate assumptions
about prime support, not full prime-power divisibility.

In a cyclic group of order $N$ with $a=v_\ell(N)$, missing $\ell$ from
an element's order is equivalent to being an $\ell^a$-th power, not merely
an $\ell$-th power when $a>1$. A uniform element misses it with
probability $\ell^{-a}$; independence and uniformity for the actual
conjugate values have not been established.
\end{definition}

\begin{theorem}[Popovych burden]\label{thm:burden}
Let $n$ be produced by the Frobenius-matching mechanism for the pair:
for every $p \mid n$ there is $j_p$ with $n \equiv p^{j_p} \pmod{
\lcm(r, \rho_r(p), \rho'_r(p))}$, and $\ord_r(n) > 2$ (validity).
Assume Definition~\ref{def:sharp} for both sections at a common
$B\ge\max_{p\mid n}2rd_p$. In the design requiring all matching orders
to be $B$-smooth, for every $p\mid n$,
writing $d_p = \ord_r(p)$ (and noting validity forces $d_p \ge 3$),
the quantity $\prod_{e \mid d_p} \Phi_e(p)$ --- a product of cyclotomic
torus orders of total dimension $\sum_{e \mid d_p}\varphi(e) = d_p$ ---
must be $B$-smooth. Prime support alone does not imply that this whole product
divides the common modulus with its full multiplicities. In
particular the minimal per-prime burden over admissible cells is
\[
\mathrm{burden}_{\mathrm{Pop}}(r) \;=\; \min\{\, d : d \mid r - 1,\
d \ge 3 \,\} \;\ge\; 3 \qquad \text{for every prime } r,
\]
whereas the single-section (Agrawal) mechanism has burden
$\min_{d \mid r-1,\, d \ge 3}(d/2$ if $d$ even, else $d)$, with global
optimum $2$. At $r = 5$: $\mathrm{burden}_{\mathrm{Pop}} = 4$, i.e.\
$p - 1$, $p + 1$ and $p^2 + 1$ simultaneously smooth, against the
single section's $2$.
\end{theorem}

\begin{proof}
For $X+2$, its order divides $p^{d_p}-1$, so smoothness of that product
is sufficient. Under Definition~\ref{def:sharp} every prime above $B$
in the product must also divide the section order, proving necessity.
This uses the support hypothesis, not Lemma~\ref{lem:nonunit} alone.
The containing product has dimension $d_p$ for both parities.
Minimizing $d_p$ over
$d \mid r - 1$, $d \ge 3$ gives the formula; the least such divisor is
$3$ iff $3 \mid r - 1$, and for $r = 5$ (where $r - 1 = 4$) it is $4$,
with levels $\{1, 2, 4\}$ and $\Phi_4(p) = p^2 + 1$. The Agrawal
comparison is Lemma~\ref{lem:collapse} plus the same minimization.
\end{proof}

\begin{remark}[heuristic cost at representative parameters]\label{cor:cost}
For pool primes of $b$ bits and smoothness bound $B$, the extra
quadratic level of Theorem~\ref{thm:burden} at $r = 5$ costs a Dickman
factor $\rho(2b / \log_2 B)$ beyond the double condition. At the
parameters of the Agrawal-directed harvest study in this repository
($b = 51$; $B = 1259$ for the solver base, $B' = 12590$ for the
enumerated base), this is $\rho(9.9) \approx 3.9\times10^{-11}$,
respectively $\rho(7.5) \approx 1.8\times10^{-7}$ --- reducing that
study's projected pool of $2.5\times10^5$ elements to an expected size
between $10^{-5}$ and $0.05$. This assumes suitable polynomial-value
smoothness and conditional independence. It is neither a proof of emptiness
at every feasible scale nor an exclusion of other structured harvests.
\end{remark}

\begin{center}
\begin{tabular}{@{}rccc@{}}
\toprule
$r$ & Agrawal burden & Popovych burden & $\Delta$ \\
\midrule
$5$ & $2$ & $4$ & $+2$ \\
$7$ & $3$ & $3$ & $0$ \\
$11$ & $5$ & $5$ & $0$ \\
$13$ & $2$ & $3$ & $+1$ \\
$17$ & $2$ & $4$ & $+2$ \\
$19$ & $3$ & $3$ & $0$ \\
$23$ & $11$ & $11$ & $0$ \\
$29$ & $2$ & $4$ & $+2$ \\
$31$ & $3$ & $3$ & $0$ \\
$37$ & $2$ & $3$ & $+1$ \\
$41$ & $2$ & $4$ & $+2$ \\
\bottomrule
\end{tabular}
\end{center}

The deltas vanish exactly at $r \equiv 3 \pmod 4$, where the Agrawal
burden is already $\ge 3$ and both conjectures share full depth; the
Agrawal global optimum ($2$, at any $r \equiv 1 \bmod 4$) is unavailable
to the pair at every modulus.

\begin{remark}[the constant $2$ is immaterial]\label{rem:constant}
For any integer $c \notin \{0, \pm 1\}$, the characteristic-zero argument of
Lemma~\ref{lem:nonunit} applies with the exclusions below. The conditional
calculation of Theorem~\ref{thm:burden} additionally requires its own
support-sharpness hypothesis for the pair $\{X-1,X+c\}$:
the proof needs only that at a prime $\mathfrak q \mid (\zeta + c)$ the
image of $-c$ has multiplicative order $r$, which holds at the primitive
prime divisors of $N(\zeta + c) = \Phi_r(-c) = (c^r + 1)/(c + 1)$. The
constant enters only through that norm (its primes, with
$p \mid c(c + 1)$, are the finitely many excluded $p$), the excluded
modulus $r \mid c + 1$ (for $c = 2$: our standing $r \ne 3$), and the
level-one component, the base-$(c{+}1)$ Fermat condition of
Remark~\ref{rem:asides}. What must be avoided is exactly degeneracy of
the zero $-c$: it must be a point of $\mathbb{G}_m$ avoiding the rational
torsion $\{\pm 1\}$, i.e.\ $c \notin \{0, \pm 1\}$. The exclusion
$c = 0$ is vacuous rather than subtle: $X$ is a monomial, so
$(X)^n \equiv X^n$ is an identity in every ring (the zero sits at the
origin, off the torus, and the section is a unit of finite order $r$).
Among nonzero $c$: $c = -1$ is the first section
itself, and $c = +1$ --- the one innocent-looking bad choice --- has its
zero at $-1$, fixed by inversion, so the collapse of
Lemma~\ref{lem:collapse} \emph{survives} for $X + 1$:
$(\zeta + 1)^{p^{d/2}} = \zeta^{-1} + 1 = \zeta^{-1}(\zeta + 1)$, whence
$\ord(X + 1) \mid r(p^{d/2} - 1)$ for even $d$; moreover $\zeta + 1 =
(\zeta^2 - 1)/(\zeta - 1)$ is multiplicatively dependent on the Galois
orbit of $\zeta - 1$, so $\{X - 1, X + 1\}$ adds no rank and no burden:
a Popovych-style conjecture with $c = 1$ would be no harder for the
mechanism than Agrawal's. Thus $c = 2$ is simply the smallest valid
choice ($c = -2$ works equally, with Mersenne norm $2^r - 1$).
Numerically (\texttt{code/ag5b\_constant\_check.py} of the repository):
at even $d$ the section $X + 1$ collapses in $22/22$ samples with the top
level never needed, while $c \in \{2, 3, -2, 5\}$ collapse only at the
single degenerate point $(r, p) = (13, 5)$ of Section~4.
\end{remark}


\begin{proposition}[the $X+1$ variant is mechanism-equivalent to
Agrawal's conjecture]\label{prop:xplus1}
Suppose $n \equiv p^{j} \pmod{\lcm(r, \rho_r(p))}$ and $p - 1 \mid n - 1$
(the matching congruence plus the Carmichael condition, as in the original
L--P sufficient construction). Then $T(1, n, r)$ holds modulo
$p$ outright. Consequently every counterexample to Agrawal's conjecture
produced by the mechanism also defeats the variant conjecture whose
hypothesis is the pair $\{T(-1,n,r),\, T(1,n,r)\}$: that variant, though
formally stronger-looking, is exactly as resistant as Agrawal's --- the
second hypothesis adds no constraint under these assumptions. The section
$X+2$ generally imposes separate order data, but this is not an exclusion
of all matching constructions.
\end{proposition}

\begin{proof}
At the $X = 1$ component, $(1+1)^n = 2^n \equiv 2 \pmod p$ since
$\ord_p(2) \mid p - 1 \mid n - 1$. In the $\Phi_r$ components,
$\zeta + 1 = (\zeta^2 - 1)/(\zeta - 1)$, a unit (its norm is
$\Phi_r(-1) = 1$), and both $\zeta - 1$ and $\zeta^2 - 1$ are values of
$X - 1$ at primitive $r$-th roots of unity, so their orders divide
$\rho_r(p) = \lcm_a \ord(\zeta^a - 1)$. Hence
$(\zeta^a - 1)^n = (\zeta^a - 1)^{p^j} = \zeta^{a p^j} - 1$ for
$a = 1, 2$, and
\[
(\zeta + 1)^n \;=\; \frac{(\zeta^2 - 1)^n}{(\zeta - 1)^n}
\;=\; \frac{\zeta^{2p^j} - 1}{\zeta^{p^j} - 1}
\;=\; \zeta^{p^j} + 1 \;=\; \zeta^n + 1,
\]
using $n \equiv p^j \pmod r$ in the last step. Measured
(\texttt{code/ag5b\_constant\_check.py}, check W): $48/48$ synthetic
mechanism instances pass $T(1, n, r)$, while $0/48$ pass $T(2, n, r)$.
\end{proof}

\section{Numerical verification}

Historical finite checks were made by a deterministic script
(\texttt{code/ag5\_popovych\_check.py} of the repository, with committed
JSON output), on $56$ sample primes spanning $(r, d) \in \{(5,4), (13,4),
(41,4), (7,6), (7,3), (11,5), (23,11)\}$: Frobenius matching for $X + 2$
($56/56$); the first section's collapse bound fails to absorb $X + 2$ at
even $d$ ($35/36$); the top level $e = d$ is present in $\rho'_r(p)$
($55/56$); large primes of all levels $e \mid d$ divide $\rho'_r(p)$
($47/47$, $\approx 0.28$ misses expected under the $1/\ell$ model);
$\gcd(r, 2^r + 1) = 1$ at all tested $r$. The single exception in the
second and third checks is $r = 13$, $p = 5$, where $\Phi_4(5) = 26 =
2 \cdot 13$ contains no prime outside $2r$ at all --- a vacuous top
level at one tiny prime, consistent with (indeed predicted by) the
generic model. It is nonetheless an accidental finite-field collapse and
does not contradict the characteristic-zero statement. These samples do
not establish the support hypotheses or asymptotic cost estimates.

\section{Scope}

Three boundaries. (i) \emph{Mechanism scope}: the results quantify
the Frobenius-matching family containing the original sufficient
construction; ``accidental''
solutions of the component equations outside it are not excluded, and no
way to construct them is known. (ii) \emph{Conditionality}: the burden
\emph{upper} structure (the collapse identities and support bounds) is
unconditional; the \emph{lower} bounds use generic sharpness
(Definition~\ref{def:sharp}), an Artin-type hypothesis --- a construction
could target primes whose section orders collapse further; no asymptotic
density or economic comparison with smooth-prime pools is proved here.
(iii) \emph{Direction}: the conditional support-model comparison is not a
universal algorithmic obstruction, an exclusion of every matching
construction, or a proof of either conjecture.

\begin{thebibliography}{9}\itemsep1pt\small

\bibitem{AKS} M.~Agrawal, N.~Kayal, N.~Saxena, \emph{PRIMES is in P},
Ann.\ of Math.\ \textbf{160} (2004), 781--793.

\bibitem{BP01} R.~Bhattacharjee, P.~Pandey, \emph{Primality testing},
BTech thesis, IIT Kanpur, 2001.

\bibitem{LPremarks} H.~W.~Lenstra, Jr., C.~Pomerance, \emph{Remarks on
Agrawal's conjecture}, AIM workshop notes, 2003.

\bibitem{Popovych09} R.~Popovych, \emph{A note on Agrawal conjecture},
Cryptology ePrint 2009/008.

\bibitem{Popovych12} R.~Popovych, \emph{Elements of high order in finite
fields of the form $\F_q[x]/\Phi_r(x)$}, Finite Fields Appl.\
\textbf{18} (2012), 700--710.

\bibitem{Vana09} T.~V\'a\v{n}a, \emph{Agrawal's conjecture and
Carmichael numbers}, Comenius University Bratislava, 2009.

\bibitem{HD21} A.~Hegde, P.~Devaraj, \emph{Heuristics for the
construction of counterexamples to the Agrawal conjecture}, Springer
PROMS \textbf{344}, 2021.

\end{thebibliography}

\end{document}
